\labsection{7}{Build and diagnose a regression model}{sec:lab07}

\begin{labproject}{Lab 7 regression: predict, benchmark, and diagnose}
\textbf{Dataset.} \texttt{ames\_housing\_dataset.csv}.\\
\textbf{Scenario.} You are building a first house-price prediction model. The goal is not to maximise fit on the available data, but to establish a transparent baseline, estimate performance on unseen records, and identify where the model fails.

\textbf{Learning outcomes.} You should be able to construct a leakage-safe train/test workflow, compare a model with a baseline, interpret holdout metrics, inspect residual structure, and evaluate whether a categorical predictor adds useful generalisable information.

\textbf{Core workflow.}
\begin{enumerate}
  \item \textbf{Define target and predictors.} Use \texttt{SalePrice} as the target and begin with living area, overall quality, year built, garage capacity, and basement area. Explain why each predictor is available before the sale outcome.
  \item \textbf{Handle missingness transparently.} Report the original row count, the rows removed or modified for selected variables, and the final modelling sample.
  \item \textbf{Split before fitting.} Create a reproducible 80/20 train/test split. Keep the test set untouched until model evaluation.
  \item \textbf{Establish a baseline.} Predict the training-set median sale price for every test record. Compute baseline MAE and RMSE before evaluating the regression model.
  \item \textbf{Fit multiple regression.} Fit the specified linear model on training data only. Interpret coefficient signs conditionally and avoid causal language.
  \item \textbf{Evaluate generalisation.} On the test set, report MAE, RMSE, and $R^2$, then quantify improvement over the baseline.
  \item \textbf{Diagnose failure modes.} Inspect predicted-versus-observed and residual-versus-fitted plots. Identify nonlinearity, changing variance, systematic under/overprediction, or influential extreme errors.
  \item \textbf{Inspect segment error.} Compare absolute error across sale-price bands or another meaningful subgroup. Determine whether a good average metric hides poor performance for a subset of homes.
\end{enumerate}

\textbf{Model-development extension.} Examine neighbourhood sample sizes and median prices before adding \texttt{Neighborhood} as a categorical predictor. Refit using the same split and report whether holdout error improves. Explain why any target-derived encoding must be learned from training data only.

\textbf{Required evidence.} Submit baseline and model metrics, two diagnostic plots, one subgroup error summary, and a short model card stating intended use, strongest observed limitation, and one next improvement to test.

\textbf{Professional standard.} A model is useful only relative to a baseline and an intended decision. Training $R^2$ alone is not evidence of predictive performance.

\tcbbreak
\lstinputlisting[style=dsr]{all_experiments/lab07_regression.R}
\end{labproject}
